\documentclass[10pt,a4paper]{amsart}
\usepackage[margin=19mm]{geometry}
\usepackage{amsmath,amssymb,mathtools}
\usepackage[scr=rsfs,cal=euler]{mathalfa}
\usepackage{xfrac}
\usepackage[unicode,hidelinks,bookmarks=false]{hyperref}
\newcommand{\ie}{i.e.\ }
\newcommand{\cf}{cf.\ }
\newcommand{\resp}{respectively\ }
\newcommand{\Subsection}{\subsection}
\providecommand{\nobreakdash}{\nobreak\hbox{-}}
\setlength{\emergencystretch}{2em}
\multlinegap0pt
\usepackage{amssymb}
\usepackage{mathtools}
\usepackage{xcolor}
\newtheorem{proposition}{Proposition}
\newtheorem{theorem}{Theorem}
\newtheorem{corollary}{Corollary}
\newcommand{\K}{\mathbb K}
\newcommand{\SR}[1]{\K[#1]}
\DeclareMathOperator{\pdim}{pdim}
\DeclareMathOperator{\reg}{reg}
\title{Explicit Betti Numbers for Skeletons of Chordal Clique Complexes and Their Alexander Duals}
\author{Mohammed Rafiq Namiq}
\date{Source: arXiv:2603.17776v2 (CC BY 4.0)}
\begin{document}
\maketitle

\noindent\emph{Source and licence.} Explicit Betti Numbers for Skeletons of Chordal Clique Complexes and Their Alexander Duals. Version arXiv:2603.17776v2 (CC BY 4.0), distributed by arXiv under the Creative Commons Attribution 4.0 International licence, http://creativecommons.org/licenses/by/4.0/, as recorded in the arXiv licence metadata for this exact version and shown on the abstract page https://arxiv.org/abs/2603.17776v2. Full licence notice: \texttt{licenses/arxiv-2603.17776-CC-BY-4.0.txt}.\par\smallskip

\noindent\textit{Editorial scope.} Counted unit: the article's theorem thm:3 (source0.tex lines 438--481). The article's complete original proof of it is delivered with it (source0.tex lines 483--534, one \texttt{proof} environment of 52 lines). No mathematics is written, repaired or re-derived: the statement and the proof are the article's own lines, in the article's own notation, and no retained line is substituted or re-flowed.  Retained uncounted as the source-local context the proof needs: lines 216--233; lines 304--345; lines 384--413; lines 536--549.  Nothing was omitted from any retained span, so the package carries no omission record.  No retained line contains a citation, so the wrapper carries no bibliography and no bibliography entry.  No retained line contains a figure environment, an image-inclusion command or drawing code, so no image is delivered and the package declares no asset dependency.  Editorial formatting only, in addition: an AMS \texttt{amsart} preamble that restates the article's own declarations and macros used by the retained lines, namely the environments \texttt{proposition}, \texttt{theorem}, \texttt{corollary} on separate counters, together with the standard packages those lines need.  The retained environments print in this document as Proposition 1, Theorem 1, Corollary 1 in order of appearance, whereas the article prints the same items with the numbering of its own full document.  The complete native source of the article as distributed in the arXiv e-print for this exact version is bundled unchanged as \texttt{sources/196646/source0.tex}.
	\begin{proposition}\label{prop:2}
		Let $\Delta_{\mathbf r}=\Delta_{\mathbf r}(n_1,\ldots,n_e)$ be the clique complex of a chordal graph with at least two maximal cliques. Fix a leaf order $S_1,\ldots,S_e$, put $n_m=|S_m|$, and let $r_1,\ldots,r_{e-1}$ be the corresponding separator cardinalities. For every integer $t\ge0$,
		\[
		f_{t-1}(\Delta_{\mathbf r})
		=
		\sum_{m=1}^e\binom{n_m}{t}
		-
		\sum_{m=1}^{e-1}\binom{r_m}{t}.
		\]
		Moreover,
		\[
		H_{\SR{\Delta_{\mathbf r}}}(z)
		=
		\sum_{m=1}^e\frac{1}{(1-z)^{n_m}}
		-
		\sum_{m=1}^{e-1}\frac{1}{(1-z)^{r_m}}.
		\]
	\end{proposition}

	\begin{theorem}\label{thm:1}
		Let $\Delta_{\mathbf r}=\Delta_{\mathbf r}(n_1,\ldots,n_e)$ be the clique complex of a chordal graph with at least two maximal cliques. Fix a leaf order $S_1,\ldots,S_e$ of its facets, put $n_m=|S_m|$, and let $r_1,\ldots,r_{e-1}$ be the corresponding separator cardinalities.
		\begin{enumerate}
			\item[(a)] For the $(-1)$-skeleton,
			\[
			\beta_{i,i}\bigl(\SR{(\Delta_{\mathbf r})_{(-1)}}\bigr)=\binom{N_{\mathbf r}}i
			\qquad(0\le i\le N_{\mathbf r}),
			\]
			and all other graded Betti numbers vanish.
			
			\item[(b)] For the $0$-skeleton,
			\[
			\beta_{i,i+1}\bigl(\SR{(\Delta_{\mathbf r})_{(0)}}\bigr)
			=i\binom{N_{\mathbf r}}{i+1}
			\qquad(1\le i\le N_{\mathbf r}-1),
			\]
			and all other Betti numbers in positive homological degrees vanish.
			
			\item[(c)] Let $1\le k<d$. For every $i\ge1$, the only possible nonzero Betti numbers in positive homological degrees are
			\[
			\beta_{i,i+1}\bigl(\SR{(\Delta_{\mathbf r})_{(k)}}\bigr)
			=
			-\sum_{m=1}^e\binom{N_{\mathbf r}-n_m}{i+1}
			+\sum_{m=1}^{e-1}\binom{N_{\mathbf r}-r_m}{i+1}
			\]
			and
			\[
			\beta_{i,i+k+1}\bigl(\SR{(\Delta_{\mathbf r})_{(k)}}\bigr)
			=
			\sum_{t=k+2}^{i+k+1}
			(-1)^{k-t}
			\binom{N_{\mathbf r}-t}{i+k+1-t}
			\left(
			\sum_{m=1}^e\binom{n_m}{t}
			-
			\sum_{m=1}^{e-1}\binom{r_m}{t}
			\right).
			\]
			
			\item[(d)] For $k=d$, the skeleton is the full complex. Its Betti numbers in positive homological degrees are given by the linear row formula in part~(c), and all other Betti numbers in positive homological degrees vanish.
		\end{enumerate}
	\end{theorem}

	\begin{theorem}\label{thm:2}
		Let $\Delta_{\mathbf r}=\Delta_{\mathbf r}(n_1,\ldots,n_e)$ be the clique complex of a chordal graph with at least two maximal cliques $S_1,\ldots,S_e$ in a fixed leaf order. Put $n_m=|S_m|$, let $r_1,\ldots,r_{e-1}$ be the corresponding separator cardinalities.
		Let $-1\le k\le d$. Then
		\[
		\reg\SR{(\Delta_{\mathbf r})_{(k)}}
		=
		\begin{cases}
			k+1,&-1\le k<d,\\
			1,&k=d,
		\end{cases}
		\]
		and
		\[
		\pdim\SR{(\Delta_{\mathbf r})_{(k)}}
		=
		\begin{cases}
			N_{\mathbf r}-k-1,&-1\le k\le r_{\min},\\
			N_{\mathbf r}-r_{\min}-1,&r_{\min}<k\le d.
		\end{cases}
		\]
		Consequently,
		\[
		\operatorname{depth}\SR{(\Delta_{\mathbf r})_{(k)}}
		=
		\begin{cases}
			k+1,&-1\le k\le r_{\min},\\
			r_{\min}+1,&r_{\min}<k\le d.
		\end{cases}
		\]
	\end{theorem}

	\begin{corollary}\label{cor:3}
		Under the hypotheses and notation of Theorem~\ref{thm:2}, for every $-1\le k\le d$, the following statements hold.
		\begin{enumerate}
			\item[(i)] The ring $\SR{(\Delta_{\mathbf r})_{(k)}}$ is initially Cohen--Macaulay if and only if
			\[
			k\le r_{\min}
			\qquad\text{or}\qquad
			r_{\min}=n_{\min}-1.
			\]
			Thus, if $r_{\min}<n_{\min}-1$, the initially Cohen--Macaulay skeletons are exactly those with $k\le r_{\min}$; if $r_{\min}=n_{\min}-1$, every skeleton is initially Cohen--Macaulay.
			
			\item[(ii)] The ring $\SR{(\Delta_{\mathbf r})_{(k)}}$ is Cohen--Macaulay if and only if $k\le r_{\min}$.
		\end{enumerate}
	\end{corollary}

	\begin{theorem}\label{thm:3}
		Let $\Delta_{\mathbf r}=\Delta_{\mathbf r}(n_1,\ldots,n_e)$ be the clique complex of a chordal graph with at least two maximal cliques $S_1,\ldots,S_e$ in a fixed leaf order. Put $n_m=|S_m|$, let $r_1,\ldots,r_{e-1}$ be the corresponding separator cardinalities. Let $c$ be the number of connected components of $\Delta_{\mathbf r}$. Every connected component of $\Delta_{\mathbf r}$ is contractible. For $1\le k<d$, put
		\[
		\lambda_k=
		\sum_{m=1}^{e}\binom{n_m-1}{k+1}
		-
		\sum_{m=1}^{e-1}\binom{r_m-1}{k+1},
		\qquad
		\mu=\bigl|\{m:r_m=r_{\min}\}\bigr|.
		\]
		Then the following statements hold.
		\begin{enumerate}
			\item[(i)] Each connected component of $(\Delta_{\mathbf r})_{(k)}$ is homotopy equivalent to a wedge of $k$-spheres, and the total number of spheres over all components is $\lambda_k$. Consequently,
			\[
			\widetilde H_q((\Delta_{\mathbf r})_{(k)};\mathbb Z)
			\cong
			\begin{cases}
				\mathbb Z^{c-1},&q=0,\\
				\mathbb Z^{\lambda_k},&q=k,\\
				0,&\text{otherwise}.
			\end{cases}
			\]
			In particular, the integral homology of every proper skeleton of positive dimension is torsion free.
			
			\item[(ii)] The last nonzero Betti numbers in the two rows are
			\[
			\beta_{N_{\mathbf r}-k-1,N_{\mathbf r}}
			\bigl(\SR{(\Delta_{\mathbf r})_{(k)}}\bigr)=\lambda_k,\quad \beta_{N_{\mathbf r}-r_{\min}-1,\,N_{\mathbf r}-r_{\min}}
			\bigl(\SR{(\Delta_{\mathbf r})_{(k)}}\bigr)=\mu.
			\]
			If $k\le r_{\min}$, the first of these is the unique extremal Betti number. If $k>r_{\min}$, these two Betti numbers are exactly the two extremal Betti numbers.
			
			\item[(iii)] If $k\le r_{\min}$, then $\SR{(\Delta_{\mathbf r})_{(k)}}$ is Cohen--Macaulay and
			\[
			\operatorname{type}\SR{(\Delta_{\mathbf r})_{(k)}}
			=
			\begin{cases}
				\lambda_k,&k<r_{\min},\\
				\lambda_k+\mu,&k=r_{\min}.
			\end{cases}
			\]
			For $1\le k<d$, the ring $\SR{(\Delta_{\mathbf r})_{(k)}}$ is level if and only if $k<r_{\min}$. If $d\ge1$, the $0$-skeleton is level of type $N_{\mathbf r}-1$. No proper skeleton of $\Delta_{\mathbf r}$ with nonnegative dimension is Gorenstein.
		\end{enumerate}
	\end{theorem}

	\begin{proof}
		Add the facets in leaf order. If $Q_m\ne\varnothing$, the simplex $S_{m+1}$ is attached to the preceding union along the nonempty face $Q_m$ and collapses onto that face relative to the preceding union. If $Q_m=\varnothing$, the new simplex forms a new connected component. It follows inductively that every connected component of $\Delta_{\mathbf r}$ is contractible. The relation tree has exactly $c-1$ edges with empty separator.
		
		Fix $1\le k<d$. The $k$-skeleton has the same connected components as $\Delta_{\mathbf r}$. If a component has dimension at most $k$, it is unchanged and hence contractible. Otherwise, the component is obtained from its $k$-skeleton by adjoining simplices of dimension at least $k+1$. Therefore the inclusion of the $k$-skeleton induces isomorphisms on homotopy groups below degree $k$. For $k\ge2$, each component of the skeleton is consequently $(k-1)$-connected; for $k=1$, each component is a connected graph. A connected graph is a wedge of circles. For $k\ge2$, the Hurewicz theorem identifies $\pi_k$ with the free abelian group $H_k$, and maps from $k$-spheres representing a basis of $H_k$ give a homology equivalence from a wedge of $k$-spheres to the component. Since both spaces are simply connected CW complexes, Whitehead's theorem makes this a homotopy equivalence. Thus, in every case, a $k$-dimensional component is homotopy equivalent to a wedge of $k$-spheres.
		
		It remains to count the spheres. Proposition~\ref{prop:2}, together with
		\[
		-1+\sum_{t=1}^{k+1}(-1)^{t-1}\binom At
		=(-1)^k\binom{A-1}{k+1}
		\qquad(A\ge1),
		\]
		gives
		\[
		\widetilde\chi((\Delta_{\mathbf r})_{(k)})
		=(c-1)+(-1)^k\lambda_k.
		\]
		Here each empty separator contributes to the term $c-1$, while under our binomial convention its contribution to $\lambda_k$ is zero. Since the only possible reduced homology groups are in degrees $0$ and $k$, the total rank in degree $k$ is $\lambda_k$. This proves part~(i).
		
		Hochster's formula on the full vertex set gives
		\[
		\beta_{N_{\mathbf r}-k-1,N_{\mathbf r}}
		\bigl(\SR{(\Delta_{\mathbf r})_{(k)}}\bigr)
		=\dim_{\K}\widetilde H_k((\Delta_{\mathbf r})_{(k)};\K)
		=\lambda_k.
		\]
		For the linear row, substitute $i=N_{\mathbf r}-r_{\min}-1$ into Theorem~\ref{thm:1}(c). Every facet term vanishes because $n_m\ge r_{\min}+1$, and a separator contributes one exactly when $r_m=r_{\min}$. Thus the second endpoint is $\mu$.
		
		Theorem~\ref{thm:1} has only the linear row and the row $j=i+k+1$. The latter ends at homological degree $N_{\mathbf r}-k-1$, while the former ends at $N_{\mathbf r}-r_{\min}-1$. Hence, when $k\le r_{\min}$, the last nonzero entry in the row $j=i+k+1$ dominates every other nonzero Betti number and is the unique extremal one. When $k>r_{\min}$, the two endpoints are incomparable and are precisely the two extremal Betti numbers. This proves part~(ii).
		
		Assume now that $k\le r_{\min}$. By Theorem~\ref{thm:2}, the projective dimension is $N_{\mathbf r}-k-1$. If $k<r_{\min}$, the linear row ends earlier, so the last free module is
		\[
		R(-N_{\mathbf r})^{\lambda_k}.
		\]
		Hence the ring is level of type $\lambda_k$. If $k=r_{\min}$, the two rows end in the same homological degree and the last free module contains
		\[
		R(-(N_{\mathbf r}-k))^{\mu}\oplus R(-N_{\mathbf r})^{\lambda_k}.
		\]
		Both summands are nonzero: $\mu>0$, and $\lambda_k>0$ because $k<d$, so there is a $(k+1)$-face whose boundary is a nonzero $k$-cycle in the $k$-skeleton. Thus the type is $\lambda_k+\mu$ and the ring is not level. When $k>r_{\min}$, Corollary~\ref{cor:3} below shows that the ring is not Cohen--Macaulay, so it is not level in the standard sense. Theorem~\ref{thm:1}(b) shows that the $0$-skeleton is level of type $N_{\mathbf r}-1$ whenever it is proper.
		
		Finally, suppose $k<r_{\min}$. Then
		\[
		\lambda_k=
		\binom{n_1-1}{k+1}
		+\sum_{m=2}^{e}
		\left[
		\binom{n_m-1}{k+1}
		-
		\binom{r_{m-1}-1}{k+1}
		\right]\ge e.
		\]
		Indeed, $S_1$ is incident with an edge of the relation tree, so $n_1\ge r_{\min}+1\ge k+2$, and for $m\ge2$ we have $n_m\ge r_{m-1}+1$ with $r_{m-1}>k$; Pascal's identity shows that every bracket is at least one. Hence the type is at least two. If $k=r_{\min}$, then $\lambda_k+\mu\ge2$. A proper $0$-skeleton has $N_{\mathbf r}\ge3$ and type $N_{\mathbf r}-1\ge2$. Therefore no proper skeleton with nonnegative dimension is Gorenstein.
	\end{proof}

\end{document}
